\section{Experimental Protocol}
\label{sec:protocol}

\paragraph*{Pre-registration}
Each comparison in this paper has a rule that was written down, with its metric, windows, floor, bound and verdict labels, in the project's decision log (\texttt{DECISIONS.md}) and committed to the repository \emph{before} the numbers it judges were committed.
% src: A16, R6
Table~\ref{tab:prereg} lists the eight pre-registrations, what each fixed, and the commits that establish the order.
The order is evidenced by commit hashes and commit-log times only (all times $+0530$, 2026-10-13 to 2026-10-15).
% src: A16 as amended by the S12 addendum: commit hashes and commit-log times only; file-timestamp and script-log times are not used.

\begin{table*}[t]
\centering
\caption{Pre-registrations. ``Rule'' is the commit that introduced the decision entry; ``numbers'' is the first commit that records a number the rule judges.}
\label{tab:prereg}
\small
\begin{tabular}{@{}lp{6.6cm}p{3.3cm}p{4.4cm}@{}}
\toprule
Entry & What it fixed & Rule (commit, time) & Numbers (commit, time) \\
\midrule
% src: git log (commit hashes and commit-log times), DECISIONS.md D43, D54, D56, D57, D59, D65, D68, D72, D76; A16
D43 & H2 conditions, the diagnostic-AUC metric over steps $[450,650)$, the verdict rule (WIN needs, among further conditions, a paired gain $\ge 0.02$ at $\rho=1$; NEGATIVE needs $\le -0.02$; NULL rules out $0.05$), seeds 1--5, draws 100--149 & 291a764, 10-13 10:57:40 (before any training) & 09cc126 (F16), 10-14 00:02:28 \\
D54 & Layer-2 attribution protocol: units, perturbation size $0.3$, floors $0.20$ and $0.05$, stability threshold $0.80$, the neighbour-clause rule, a $100$\,ms timing assumption & aa97280, 10-14 15:57:15 & 3487836 (F24), 10-14 16:15:56 \\
D56 + D57 & the sweeps, the H1 and H3 metrics, the comparisons and the WIN/NULL/NOISY rule; D57 only re-normalised the H1 effect by A's count, because the oracle's count is zero by construction, and the entry records that it was written before any evaluation draw produced a result & 26e6e79 (D56), 10-14 16:53:33; 9f2a8f9 (D57), 10-14 17:05:45 & 0b6585e (F26), 10-14 17:59:53 \\
D59 & Gate-C: the $\tau$ grid and its tuning, the evaluation points, the ``weighted beats gate'' criterion & 3710795, 10-14 19:26:16 & 86b3a20 (first Gate-C code), 19:29:33; eead7f1 (F28, tuning grid), 19:29:55; ea44011 (F29, results), 19:37:51 \\
D65 & monitor M (weights by the D52 procedure; verdict rule), staggered onset ($S=300$), the mechanism check that gates Experiment 2, the latency rule with a $\pm 20$-step band & 245705b, 10-15 09:16:12 & 776beff (F32, M tuning grid), 09:21:52; b2c09c7 (F35, mechanism check), 12:06:23 \\
D68 & the contained head D2 (architecture, parameter counts), draws 250--299, the primary latency rule ($\pm 20$ steps), no further architecture experiment & ac72b82, 10-15 13:34:48 & fd167d5 (F43, tuning grids), 14:37:34; e74a1a0 (F44, results), 17:22:46 \\
D72 & the forward-looking objective (a 150-step horizon, threshold $0.75$, minimum 5 tasks), C-h and D2-h with no architecture change, draws 300--349, the primary latency rule ($\pm 20$ steps), the last graph experiment & e7c7d54, 10-15 18:51:15 & 6890d59 (F47, tuning grids), 19:03:03; da1b039 (F48, results), 19:08:09 \\
D76 & the confirmatory test of C-h against C: fresh draws 350--399, all degraded RSUs, co-primary latency ($\pm 20$ steps) and false-alarm rate (bound $0.05$), the joint verdict, no training or tuning & 96246a8, 10-15 22:24:12 & f2eb79d (F50, results), 22:42:31 \\
\bottomrule
\end{tabular}
\end{table*}

D57 was a correction made before the evaluation, not a change of rule after it: the harness smoke test (reserved debugging draws, not results) had exposed a zero denominator in the H1 formula, and the amendment changes only that formula; every floor, bound and window of D56 stands.
% src: A16 (D57 = single amendment made before any evaluation draw ran), X14 (describe D57 only as a pre-evaluation correction); D57 text
The pre-registrations of the staggered-onset, contained-head and horizon-objective experiments additionally fix a single pre-declared extension if a primary verdict is NOISY; it was not triggered in any of them. The confirmatory test allows no extension.
% src: D65 section 2.5, D68 section 4, D72 section 4, D76 section 6; A25 (extension not triggered), F45, F49

\paragraph*{Seeds, draws, labels}
Seeds 1--5 with frozen checkpoints; evaluation draws 100--149 (the sweeps), 200--249 (staggered onset), 250--299 (D2), 300--349 (horizon objective) and 350--399 (the confirmatory test); training draws are below 100; the tuning seed (100) and the demo draw (seed 100, draw 7) are excluded from every evaluation.
% src: A16, A25, A28, A29, A30
The unit of analysis is the seed: per-seed means of per-draw paired differences, with 95\% Student-$t$ intervals over 5 seeds (4 degrees of freedom).
We use no $p$-values and no multiplicity correction; ``the interval excludes $0$'' and ``lies above the floor'' are the only inferential statements.
% src: A16, X11
Draws within a seed share a model and traffic and are not independent samples, so no draw-level standard error is reported.
Every sweep point reports 250 draws attempted, 0 failed and 250 aggregated, except where a section states otherwise (draws with no late-onset RSU are not aggregated in the staggered latency primary).
% src: A16 (D42 accounting: 250 attempted, 0 failed, 250 aggregated); A25, A28 (164 / 183 aggregated at rho 0)

\paragraph*{Verdict labels}
The rules share one form, with the exact conditions in the cited entries: with a floor $f$ and a bound $b$ fixed in advance, an effect with interval $[lo,hi]$ is WIN if $lo>0$ and the mean is at least $f$; NEGATIVE if $hi<0$ and the mean is at most $-f$; NULL if neither holds and the interval rules out an effect of size $b$; and NOISY otherwise.
% src: D43 section 4, D56 section 4, D59 section 4, D65 sections 1.4 and 2.5, D68 section 4 (the shared form; floors differ per metric)
The floors are $0.02$ (AUC), $0.01$ (success), $0.10$ of A's count (dispatches) and $20$ steps (latency), and are stated assumptions.
% src: D43 (0.02), D56 (H3 floor 0.01, H1 floor 0.10), D65/D68 (20 steps)
The horizon-objective experiment uses the same latency rule. The confirmatory test names its latency labels EARLIER and LATER (for WIN and NEGATIVE, floor $20$ steps) and labels the false-alarm difference ACCEPTABLE only if the upper end of its interval is at most $0.05$; its verdict is OBJECTIVE HELPS only if both hold, and NOT CONFIRMED otherwise.
% src: D72 section 4; D76 sections 4-6; A30
We quote each label together with the signed difference and its interval, and a label that is NULL under a floor is not a statement that the difference is zero.
% src: R4, X18 (the NULL label states a floor, not a sign)

\paragraph*{Disclosure}
The run of the contained-head experiment was interrupted once after the closed-loop runs and raw tables had been produced and before their raw results were committed; the results were regenerated from the stored files by the same commands, and nothing was re-run or re-tuned.
% src: D69 item 6 (process disclosure). NOT an A-entry: see S12 report.

\paragraph*{Detection latency}
% src: CLAIMS D78 caveat (A7, A9, A23, A25, A28-A30), stated once here and referred to elsewhere (presentation choice, see DECISIONS); A19 (false-alarm sentence); F51
Every detection latency in this paper is compared under one shared rule, not at a detection operating point.
The rule's false-alarm rate was 0.32--0.44 of healthy RSUs where it was measured (the confirmatory test, heads C and C-h only) and was not measured elsewhere.
A latency difference is therefore a comparison under this rule, not the delay of a detector run at a chosen false-alarm rate; wherever a latency appears below, this caveat applies.

\paragraph*{Code, data and decision log}
% src: A17 (reproducible from config and seed; figures and tables regenerate from one command, byte-identical on re-run); A16 (commit order)
The simulator, models, experiment scripts, the decision log (\texttt{DECISIONS.md}), the findings log (\texttt{FINDINGS.md}) and the full commit history are public at \url{https://github.com/mucool-codes/TrustGraph}, where the commits of Table~\ref{tab:prereg} can be checked; the version used for this paper is tagged \texttt{s12d-complete}.
Every run is reproducible from a configuration file and a seed, and every figure and table in this paper is regenerated by one script.
